\section{The Fibonacci generator renewal}\label{sec:fibonacci-renewal}

Let \(F_0=0,F_1=1\), and \(F_{j+2}=F_{j+1}+F_j\).  The classical Fibonacci
partition function is
\begin{equation}
 R(N):=\card{\left\{(\varepsilon_j)_{j\ge0}\in\{0,1\}^{(\mathbb N)}:
 N=\sum_{j\ge0}\varepsilon_jF_{j+2}\right\}},
 \qquad N\ge0.
 \label{eq:fibonacci-partition-function}
\end{equation}
Thus the available parts are \(1,2,3,5,\ldots\), each used at most once.
For \(m\ge1\), put
\begin{equation}
 \mathcal I_m=[F_{m+1}-1,F_{m+2}-1)\cap\ZZ,
 \qquad
 Z_m^R(t):=\sum_{N\in\mathcal I_m}R(N)^t.
 \label{eq:standard-fibonacci-layer}
\end{equation}

We use the generator coordinates of Weinstein
\cite[Secs.~3--5]{Weinstein2016}.  His Fibonacci sequence
\(f_0=f_1=1\) satisfies \(f_j=F_{j+1}\), so his half-open layer
\([f_m-1,f_{m+1}-1)\) is exactly \(\mathcal I_m\), and his partition
function is \(R\) in \eqref{eq:fibonacci-partition-function}.  We record the
part of the parameterization needed for weighted sums.

Let
\[
 \mathcal A=\{p/q:0<p<q,\ (p,q)=1\}.
\]
For \(p/q\in\mathcal A\), write its unique negative continued fraction as
\[
 \frac pq=
 \cfrac1{e_1-\cfrac1{e_2-\ddots-\cfrac1{e_k}}},
 \qquad e_i\ge2,
\]
and recall from \Cref{sec:brocot-condensation} that
\[
 d(p/q)=\sum_{i=1}^k(e_i-1),
 \qquad c(p/q)=2d(p/q)+1.
\]
For real \(s\), define
\begin{equation}
 b_\ell(s)=\sum_{\substack{p/q\in\mathcal A\\c(p/q)=\ell}}q^{-s},
 \qquad
 B_s(z)=\sum_{\ell\ge1}b_\ell(s)z^\ell.
 \label{eq:renewal-letter-series}
\end{equation}
Only odd costs occur.  The smallest costs are \(3\) and \(5\), represented
by \(1/2\) and \(1/3\).

Put
\begin{equation}
 u_0(s)=1,
 \qquad
 u_j(s)=\sum_{\ell=1}^jb_\ell(s)u_{j-\ell}(s)\quad(j\ge1).
 \label{eq:weighted-renewal-coefficients}
\end{equation}
Equivalently, in formal power series,
\begin{equation}
 U_s(z):=\sum_{j\ge0}u_j(s)z^j=\frac1{1-B_s(z)}.
 \label{eq:weighted-renewal-generating-function}
\end{equation}

\begin{proposition}[Weighted generator renewal]
\label{prop:weighted-generator-renewal}
For every real \(s\) and \(m\ge1\),
\begin{equation}
 Z_m^R(-s)
 =1+u_m(s)+2\sum_{j=1}^{m-1}u_j(s)
 =2\sum_{j=0}^{m-1}u_j(s)+u_m(s)-1.
 \label{eq:standard-layer-renewal}
\end{equation}
Moreover, \(u_j(s)\) is exactly the total inverse-\(s\) weight of the
Weinstein generators of cost \(j\).
\end{proposition}

\begin{proof}
Under Weinstein's free-monoid bijection, a generating number \(g\)
corresponds uniquely to a nonempty word
\[
 w=\frac{p_1}{q_1}\times\cdots\times\frac{p_r}{q_r}
 \qquad(p_i/q_i\in\mathcal A).
\]
Proposition~3.3 and formula~(16) of
\cite{Weinstein2016}, translated through the preceding Fibonacci convention,
give
\begin{equation}
 R(g)=\prod_{i=1}^rq_i,
 \qquad
 L(g)=\sum_{i=1}^rc(p_i/q_i),
 \label{eq:generator-weight-cost}
\end{equation}
where \(L(g)\) is the largest Zeckendorf index of \(g\).  Unique
factorization in the free monoid and \eqref{eq:renewal-letter-series} show
that the total weight \(R(g)^{-s}\) of generators with \(L(g)=j\) is the
coefficient \(u_j(s)\).  They also prove
\eqref{eq:weighted-renewal-generating-function}.

It remains to count one generating orbit in a fixed half-open Fibonacci
layer.  Weinstein's normal form consists of the four branches
\([a]g,[a]\tau g,\sigma([a]g),\sigma([a]\tau g)\), with \(a\ge0\).  His
formulas~(17)--(19) give the corresponding largest indices
\[
 \begin{array}{c|c}
 \text{branch}&\text{largest index}\ \\ \hline
 {}[a]g&L(g)+2a\\
 {}[a]\tau g&L(g)+1+2a\\
 \sigma({}[a]g)&L(g)+1+2a\\
 \sigma({}[a]\tau g)&L(g)+2+2a.
 \end{array}
\]
The two middle entries are distinct because the action is free.  Hence an
orbit of generator cost \(j=L(g)\) contributes one point to
\(\mathcal I_j\), two points to every \(\mathcal I_m\) with \(m>j\), and
no point to earlier layers.  Distinct generators have disjoint orbits.
There is in addition exactly one partition of multiplicity \(1\) in each
\(\mathcal I_m\), namely its left endpoint.  Summing the orbit weights
therefore gives the first expression in \eqref{eq:standard-layer-renewal};
the second is algebraic rearrangement.
\end{proof}

The critical normalization follows directly from the totient Dirichlet
series.  For \(s>2\), every denominator \(q\ge2\) occurs with its
\(\varphi(q)\) reduced numerators, so Tonelli's theorem gives
\begin{equation}
 B_s(1)=\sum_{q\ge2}\frac{\varphi(q)}{q^s}
 =\frac{\zeta(s-1)}{\zeta(s)}-1.
 \label{eq:weighted-renewal-at-one}
\end{equation}
The right side decreases continuously from infinity to zero as \(s\) runs
from \(2\) to infinity.  Thus there is a unique \(\sigma_0>2\) with
\begin{equation}
 \frac{\zeta(\sigma_0-1)}{\zeta(\sigma_0)}=2,
 \qquad B_{\sigma_0}(1)=1.
 \label{eq:critical-renewal-normalization}
\end{equation}

At this exponent, \(b_j(\sigma_0)\) is an aperiodic probability law: its
support contains \(3\) and \(5\).  Its mean
\begin{equation}
 \mu_C:=\sum_{q\ge2}\sum_{\substack{1\le p<q\\(p,q)=1}}
 c(p/q)q^{-\sigma_0}
 \label{eq:critical-cost-mean}
\end{equation}
is finite.  Indeed, conversion to the canonical regular continued fraction
\([0;a_1,\ldots,a_r]\) gives
\(c(p/q)=2(a_1+\cdots+a_r)-1\).  The fixed-denominator mean theorem of
Panov and Liehl \cite{Panov1982,Liehl1983} bounds the sum of
\(a_1+\cdots+a_r\) over reduced numerators by
\(O(\varphi(q)(\log q)^2)\).  Therefore
\[
 \mu_C\ll\sum_{q\ge2}q^{1-\sigma_0}(\log q)^2<\infty.
\]
This verifies the finite-mean hypothesis used in the critical renewal
argument below.
